\documentclass[11pt]{article}
\usepackage[margin=1in]{geometry}
\usepackage{amsmath,hyperref}
\setlength{\emergencystretch}{1em}
\title{Build Record and Corrections to the 2021 UAV Design}
\author{Liam Cawley}
\date{October 2026}
\begin{document}
\maketitle

The archived design report follows this note. It combines proposed mission
requirements, component choices, and engineering estimates. Several statements
need correction before the document is used as a record of the completed build.

\section*{The aircraft that was built}
The completed airframe was the 1.6 m MFD Crosswind Mini. It integrated a Pixhawk
Cube Orange running ArduPilot, RTK GPS, telemetry, and a camera payload.
Autonomous takeoff and waypoint flights were conducted. The later vision
pipeline used YOLOv5 with Sony A6000 imagery.
The report's 1.9 m Crosswind/Nimbus Pro and Keras letter-classifier descriptions
belong to the earlier design. Their dimensions and balance target should not
be transferred to the Mini as measured specifications.

\section*{Competition requirements and certification}
The four-mile course, 250,000 m$^2$ imaging task, and 40-minute mission budget
were competition requirements. They are not measured endurance or mapping
performance in this report. Completing an AUVSI SUAS mission does not confer a
federal ``level-5 autonomy'' certification. The original conclusion conflated a
competition demonstration, an autonomy classification, and regulatory approval.

\section*{Electrical estimates}
Motor $K_v$ is a speed constant in rpm per volt, not kilovoltage. The current
calculation in the report uses a 4300 $K_v$, 4.1-inch propeller example even
though the selected system is described as 800 $K_v$ with an 8-by-6 propeller.
That calculation does not establish the selected propulsion system's current.
Its 30 A cruise and 65 A acceleration figures should not be presented as
measurements.

There is also a useful consistency check: at 60 A total current, a 6 Ah battery
provides an ideal $6/60=0.1$ hours, or six minutes, before accounting for reserve
and losses. The quoted currents cannot support a 40-minute endurance claim.
Endurance needs measurements from the actual battery, propellers, aircraft
weight, and flight profile.

\section*{Control and geolocation}
A PID controller produces corrective control commands from feedback; it does
not create a gyroscopic effect or eliminate all disturbances. Trim balance
alone also does not establish static stability or stall margins.

Target geolocation from an image requires camera geometry and aircraft pose
and altitude, together with an appropriate ground model. Elapsed time and
distance along the flight path alone do not determine the coordinates of an
arbitrary pixel. No localization-error measurement is supplied by the report.

\section*{Unverified mission components}
The parachute drop, ground-vehicle survival, self-righting behavior, and
autonomous ground traversal are design requirements in the document. A wide
wheelbase alone does not establish self-righting, and the report includes no
drop-test record supporting survival from 150 ft if the parachute fails.
Agricultural cost or emissions savings likewise require a separate measured
comparison; the motivating discussion does not establish them.

The retained photographs document bench work and airframe integration. They
complement the autonomous-flight account, without supplying the mission-level
performance measurements listed above.
\end{document}
